\documentclass[preprint,12pt,authoryear]{elsarticle}

\usepackage{amsmath,amssymb,mathtools}
\usepackage{booktabs,multirow,array}
\usepackage{graphicx}
\usepackage{xcolor}
\usepackage{enumitem}
\usepackage{lineno}
\usepackage{microtype}
\usepackage{url}

\journal{Computers \& Industrial Engineering}
\modulolinenumbers[5]

\definecolor{tbdred}{RGB}{170,0,0}
\newcommand{\TBD}[1]{\textcolor{tbdred}{[TBD: #1]}}
\newcommand{\method}{SA-RLCS}
\newcommand{\cmax}{C_{\max}}
\newcommand{\pdi}{\mathrm{PDI}}
\newcommand{\gap}{\mathrm{Gap}}

\begin{document}

\begin{frontmatter}

\title{Structure-Aware Reinforcement Learning for Cutting-Plane Selection in Dual-Source Mixed-Flow Semiconductor Scheduling}

\author[inst1]{\TBD{First author}}
\author[inst1]{\TBD{Co-author}}
\author[inst2]{\TBD{Corresponding author}\corref{cor1}}
\ead{\TBD{corresponding.author@email}}
\cortext[cor1]{Corresponding author.}

\affiliation[inst1]{organization={\TBD{Department and University}},
                    city={\TBD{City}},
                    postcode={\TBD{Postcode}},
                    country={\TBD{Country}}}
\affiliation[inst2]{organization={\TBD{Department and University}},
                    city={\TBD{City}},
                    postcode={\TBD{Postcode}},
                    country={\TBD{Country}}}

\begin{abstract}
Dual-source mixed-flow semiconductor cluster tools couple robotic transport,
load locks, shared rotating chambers, batching, reusable process-environment
carriers, cleaning rules, and time-residency requirements. Their exact
mixed-integer programming (MIP) models are consequently difficult to solve
within operational time limits. This study develops a structure-aware
reinforcement-learning cut selector (\method{}) for an exact compact
formulation of this scheduling problem. The selector augments 13 generic cut
descriptors with 10 formulation-role features that quantify how a cut couples
binary decisions, general integers, objective variables, and physical
auxiliaries. A hierarchical policy chooses a cardinality and an ordered set of
policy anchors; a monotone role-submodular completion rule then balances cut
quality, structural coverage, and representativeness. The callback only ranks
solver-generated valid inequalities and returns every omitted cut in legal
order, so it does not alter the feasible region or the exactness of
branch-and-cut. Training minimizes a primal--dual integral aligned with
anytime industrial performance. A frozen paired protocol compares solver
defaults, hierarchical baselines, feature and search ablations, and the full
method on nominal, factorial, sensitivity, selective pre-batching, and
out-of-distribution instances. \TBD{Insert the validated main numerical
result, statistical conclusion, and OOD result after all frozen campaigns
finish.} The proposed framework provides a reproducible route for combining
manufacturing structure with learning-assisted exact optimization.
\end{abstract}

\begin{highlights}
\item An exact compact model preserves mixed-flow cluster-tool physics.
\item A 23-dimensional role-aware representation characterizes candidate cuts.
\item Policy anchors and submodular completion balance quality and diversity.
\item The learned callback selects only solver-generated valid inequalities.
\item A frozen paired protocol covers nominal, factorial, and OOD regimes.
\end{highlights}

\begin{keyword}
semiconductor manufacturing \sep cluster tool scheduling \sep
mixed-integer programming \sep cutting-plane selection \sep
reinforcement learning \sep submodular optimization
\end{keyword}

\end{frontmatter}

\linenumbers

\section{Introduction}
\label{sec:introduction}

Cluster tools integrate several processing modules and robotic handlers around
a tightly coupled transport system. They reduce footprint and contamination,
but a schedule must coordinate processing, transfers, chambers, load locks,
carrier availability, and recipe-dependent temporal restrictions. The
difficulty is amplified in a dual-source mixed-flow setting: product wafers
and reusable process-environment carriers enter through different sources,
meet inside shared chambers, and follow different but interacting routes.
Batch and paired processing modes further create discrete choices whose
effects propagate through the entire schedule.

Exact MIP is attractive because it can certify schedule quality and encode
physical constraints transparently. Its practical limitation is the large
branch-and-cut search induced by sequencing and resource-conflict
disjunctions. General-purpose solvers generate many candidate cuts at a node,
and adding all of them can strengthen the relaxation while also increasing
linear-program size, numerical burden, and callback overhead. Cut selection
is therefore a sequential decision problem in which local efficacy alone does
not capture interactions among cuts.

Recent work has shown that machine learning can support combinatorial
optimization without replacing the underlying mathematical model
\citep{bengio2021machine}. Reinforcement learning has been applied to cut
selection \citep{tang2020cut}, look-ahead labels
\citep{paulus2022looking}, and hierarchical sequence decisions
\citep{wang2023hem}. However, generic cut features cannot explicitly identify
whether a candidate mostly concerns manufacturing choices, timing
auxiliaries, or the makespan. A policy trained only on generic statistics may
therefore select individually attractive but structurally redundant cuts.

This paper studies the following research question:
\emph{Can formulation roles be incorporated into a learned cut selector to
improve the anytime behavior of an exact cluster-tool scheduling model while
preserving solver validity and reproducibility?} We answer the question with
an exact compact model and a structure-aware selector, denoted \method{}. The
model removes variables fixed by manufacturing equalities through continuous
relaxation while preserving the same feasible schedules; genuine assignment
and ordering choices remain integral. The selector uses a 23-dimensional
representation, a hierarchical policy, and a role-submodular completion
operator.

The paper makes four contributions.

\begin{enumerate}[leftmargin=*]
  \item We formulate a dual-source mixed-flow cluster-tool scheduling problem
  with shared robotic resources, two chamber modes, early cleaning,
  carrier-residency constraints, and reusable carrier tokens in an exact
  compact MIP.
  \item We introduce 10 formulation-role descriptors in addition to 13
  generic cut descriptors, providing a portable representation of how a cut
  couples discrete decisions, objective variables, and continuous physical
  auxiliaries.
  \item We combine learned policy anchors with a monotone submodular
  completion rule. The construction explicitly rewards quality, role
  coverage, and representativeness while enforcing a fixed cut budget.
  \item We define a frozen, instance-paired evaluation protocol covering
  nominal comparison, controlled ablations, factorial conditions, selective
  pre-batching, one-factor sensitivity, and size/recipe distribution shifts.
\end{enumerate}

The current document is an evidence-disciplined draft: numerical cells remain
blank until every required paired campaign is complete. This prevents
incomplete seeds or successfully solved instances from being reported
selectively.

\section{Related work}
\label{sec:related}

\subsection{Semiconductor cluster-tool scheduling}

Cluster-tool scheduling combines machine scheduling with robotic transport and
residency constraints. Exact and analytical methods have been developed for
photolithography areas containing cluster tools
\citep{madathil2018photolithography} and for time-constrained tools with purge
operations \citep{yuan2023purge}. These studies demonstrate the value of
representing tool physics explicitly. The present setting differs by
simultaneously modeling two material sources, product and carrier flows,
shared rotating chambers, two recipe modes, cleaning decisions, and reusable
carrier tokens.

\subsection{Learning within exact MIP}

SCIP provides a modular branch-cut-and-price framework in which learning can
act at controlled decision points \citep{achterberg2009scip,bestuzheva2023scip}.
Learning-assisted exact optimization is most defensible when the learned
component changes search priorities rather than mathematical validity. This
principle motivates our callback: it receives only valid cuts generated by
SCIP, selects a preferred subset, and leaves omitted cuts available to the
solver. The learned model neither creates inequalities nor prunes nodes.

\subsection{Learned cut selection}

\citet{tang2020cut} formulate cut choice as reinforcement learning, while
\citet{paulus2022looking} use look-ahead information to learn which cuts yield
useful bound improvement. \citet{wang2023hem} model the number and order of
selected cuts hierarchically, drawing on pointer-network ideas
\citep{vinyals2015pointer}. We retain the hierarchical count/sequence
principle but add formulation-role features and deterministic submodular
completion. The distinction is important: a neural score can express a
learned prior, whereas coverage and diversity have transparent diminishing
returns and can be imposed directly.

\section{Problem description and exact compact formulation}
\label{sec:model}

\subsection{Manufacturing topology}

Product wafers follow the route
\[
 \mathrm{LP}\rightarrow\mathrm{ATR}\rightarrow\mathrm{AL}
 \rightarrow\mathrm{ATR}\rightarrow\mathrm{LL_U}
 \rightarrow\mathrm{VTR}\rightarrow\mathrm{CH}
 \rightarrow\mathrm{VTR}\rightarrow\mathrm{LL_L}
 \rightarrow\mathrm{ATR}\rightarrow\mathrm{LP},
\]
where LP is the load port, ATR and VTR are atmospheric and vacuum transfer
robots, AL is the aligner, and $\mathrm{LL_U}$ and $\mathrm{LL_L}$ are the
upper and lower load locks. Reusable process-environment carriers (PECs)
travel between PEC storage and a shared chamber through the VTR. Chambers
CH2 and CH3 rotate and share transfer/load resources.

Two recipe modes are considered. In the $4\times1$ mode, four product wafers
are processed as a batch using one of two reachable slot groups. In the
$2\times2$ mode, a PEC and two product wafers form a physically ordered
PEC--product--product--PEC chain. Cleaning can be triggered early subject to
the cleaning interval; the selected PEC must remain resident for its required
chain and is represented by a reusable token. Resource locks prevent
simultaneous robot moves and incompatible loading operations.

\subsection{Event and resource variables}

Let $\mathcal{O}$ be the operations, $\mathcal{R}$ the physical resources,
and $\mathcal{A}$ the precedence arcs induced by routes. Each operation
$i\in\mathcal{O}$ has start time $s_i$, processing duration $p_i$, and
completion time $c_i=s_i+p_i$. Binary variables $y_{ijr}$ orient two
operations competing for unit-capacity resource $r$. Assignment and
mode-selection variables are denoted by $x$, and PEC-token assignments by
$u$. The model includes the precedence constraints
\begin{equation}
  s_j \ge c_i,\qquad (i,j)\in\mathcal{A},
  \label{eq:precedence}
\end{equation}
and resource disjunctions
\begin{align}
  s_j &\ge c_i-M_{ijr}(1-y_{ijr}), \\
  s_i &\ge c_j-M_{jir}y_{ijr},
  \qquad \{i,j\}\subseteq\mathcal{O}_r,
  \label{eq:disjunction}
\end{align}
with event-specific valid constants $M_{ijr}$. Capacity, chamber-rotation,
slot-group, cleaning, residency, and token-conservation constraints link
$x$, $u$, $y$, and $s$. The complete index-level formulation and parameter
construction are provided in \TBD{the online supplement and repository link}.

\subsection{Exact compactness}

Some nominally binary membership, slot, and alias variables are already fixed
to zero or one by equalities induced by a selected manufacturing mode. We
declare these variables continuous in the compact formulation while retaining
the fixing equalities. This operation changes neither their values nor the
projection of the feasible region; it only removes unnecessary integrality
declarations. All genuine assignments, orderings, mode choices, and cleaning
choices remain integer. Accordingly, ``compact'' denotes an equivalent
reformulation rather than a relaxation of manufacturing physics.

\subsection{Lexicographic objective}

The primary objective is the interval from the first product departure at LP
to the last product return:
\begin{equation}
  \min \cmax .
  \label{eq:makespan}
\end{equation}
For equal $\cmax$, the secondary objective reduces avoidable imbalance. If
resource $r$ has capacity $k_r$ and fixed workload $W_r$, its total idle time
is
\begin{equation}
  I_r=k_r\cmax-W_r.
\end{equation}
We minimize the lexicographic pair
\begin{equation}
 \min_{\mathrm{lex}}\left(
   \cmax,\;
   \sum_{r\in\mathcal{R}_{\mathrm{key}}} I_r^2
   +\eta\sum_{g\in\mathcal{G}} w_g
 \right),
 \label{eq:lexicographic}
\end{equation}
where the five key physical resource groups are fixed before evaluation and
$w_g$ measures avoidable waiting. The small linear term only breaks ties
remaining after makespan and squared idle time.

\section{Structure-aware learned cut selection}
\label{sec:method}

\subsection{Decision interface and learning objective}

At a separation round, SCIP supplies a legal candidate pool
$\mathcal{P}=\{1,\ldots,n\}$. Candidate $i$ represents a valid inequality
$a_i^\top z\le b_i$. A low-cost efficacy prefilter retains at most 128
candidates. The policy selects no more than 16 cuts; omitted candidates are
returned in their original legal order.

The state contains candidate features and solver context. An action consists
of a selected cardinality and an ordered subset. Training uses the negative
primal--dual integral as reward,
\begin{equation}
 R=-\pdi=-\int_0^T g(t)\,dt,
 \label{eq:pdi}
\end{equation}
where $g(t)$ is the normalized primal--dual gap and $T$ is the time limit.
Lower PDI means better anytime performance. With baseline $b$, the policy
gradient estimator is
\begin{equation}
 \widehat{\nabla_\theta J}
 =\frac{1}{B}\sum_{\ell=1}^{B}
   (R_\ell-b_\ell)\nabla_\theta
   \log\pi_\theta(a_\ell\mid s_\ell),
 \label{eq:reinforce}
\end{equation}
following \citet{williams1992reinforce}. Independent training samples can be
evaluated in parallel, but the estimator is reported as policy gradient
rather than as an asynchronous actor--critic architecture.

\subsection{Generic and formulation-role features}

The generic 13-dimensional vector contains objective parallelism, efficacy,
support, integral support, normalized violation, and the mean, maximum,
minimum, and standard deviation of cut and objective coefficients.

The structural extension partitions variables into five formulation roles:
binary decisions, general integers, continuous objective variables,
continuous physical auxiliaries, and other variables. For cut $i$ and role
$k$, coefficient-mass share is
\begin{equation}
 m_{ik}=
 \frac{\sum_{j\in\mathcal{V}_k}|a_{ij}|}
      {\sum_j|a_{ij}|+\epsilon}.
 \label{eq:mass}
\end{equation}
The five shares are augmented by bounded-variable ratio, positive-coefficient
ratio, discrete--continuous coupling, dominant-role share, and normalized
role entropy. The resulting 23-dimensional vector is scale-stable and relies
on formulation roles rather than instance-specific variable names.

\subsection{Hierarchical policy and beam search}

A count policy first predicts the selected proportion. A pointer-style policy
then scores an ordered sequence \citep{vinyals2015pointer}. The network uses a
128-dimensional hidden representation and one glimpse. Training samples
stochastically; evaluation uses either greedy decoding or beam width three.
End-of-sequence handling, accumulated log probabilities, and backpointers are
shared across learned baselines to ensure that a beam-search comparison does
not change candidate pools or accounting.

\subsection{Role-submodular completion}

Let $A$ be a small set of policy anchors and $S\supseteq A$ the completed set.
For normalized quality $q_i$, policy score $p_i$, role mass $m_{ik}$, feature
embedding $h_i$, and candidate pool $\mathcal{P}$, completion maximizes
\begin{align}
F(S)=&\;w_q\sum_{i\in S}q_i+w_p\sum_{i\in S}p_i
 +w_c\sum_k \sqrt{\sum_{i\in S}m_{ik}} \nonumber\\
&+\frac{w_r}{|\mathcal{P}|}
 \sum_{j\in\mathcal{P}}\max_{i\in S}
 \operatorname{sim}(h_i,h_j),
\label{eq:submodular}
\end{align}
subject to $|S|\le K$. Default normalized weights are
$(w_q,w_p,w_c,w_r)=(0.35,0.20,0.25,0.20)$, the completion pool factor is
2.0, and the anchor ratio is 0.25.

With nonnegative modular terms, concave-over-modular role coverage, and
facility-location representativeness, $F$ is monotone submodular. Greedy
completion of the residual cardinality budget therefore inherits the
classical $(1-1/e)$ guarantee for that fixed-anchor subproblem
\citep{nemhauser1978analysis}. This is a guarantee for the set-completion
objective, not for branch-and-cut runtime or final MIP quality.

\subsection{Validity and overhead}

\method{} never generates a constraint, changes a coefficient, or declares a
node infeasible. Every selected inequality has already passed SCIP's legality
checks. Unselected cuts remain available according to the callback contract.
Thus the learned component can alter search trajectory and runtime but not
the scheduling model's feasible set or the validity of an eventual proof.
All callback inference, beam search, and postprocessing time is charged to
the common wall-clock limit.

\section{Experimental design}
\label{sec:experiments}

\subsection{Frozen benchmark suites}

Table~\ref{tab:suites} summarizes the frozen suites. Nominal instances are
used for the principal comparison. The factorial suite varies wafer count,
recipe mix, and processing-time scale. Sensitivity conditions perturb one
physical factor at a time. The OOD suite increases problem size and changes
recipe balance. No test instance is used to tune hyperparameters.

\begin{table}[htbp]
\centering
\caption{Frozen evaluation suites.}
\label{tab:suites}
\begin{tabular}{p{0.19\linewidth}p{0.51\linewidth}r}
\toprule
Suite & Design & Conditions/instances \\
\midrule
Nominal main & Frozen nominal manufacturing instances & 20 \\
Factorial DOE & Wafers $\{8,16,24,32\}$, five recipe ratios,
processing scale $\{0.9,1.0,1.1\}$ & 60 \\
SPBS & Non-pure-$4{:}1$ conditions, none vs.\ automatic
selective pre-batching & 48 pairs \\
Sensitivity & Motion $\{0.8,1.2\}$, processing $\{0.8,1.2\}$,
PEC capacity 10, cleaning interval $\{3,10\}$ & 8 \\
OOD & Wafers $\{40,48,64\}$ and ratios
$\{3{:}1,1{:}1,1{:}3\}$ & 9 \\
\bottomrule
\end{tabular}
\end{table}

\subsection{Compared methods}

\begin{table}[htbp]
\centering
\caption{Methods in the frozen main comparison. All methods use the same
compact formulation, warm start, SCIP configuration, candidate prefilter,
time accounting, and solver seeds.}
\label{tab:methods}
\begin{tabular}{p{0.22\linewidth}p{0.69\linewidth}}
\toprule
Method & Description \\
\midrule
SCIP & Default SCIP cut selection. \\
ACS & Deterministic adaptive cut-selection baseline tuned only on the frozen
validation set. \\
HEM & 13 generic features, hierarchical policy, greedy decoding, no
role-aware completion. \\
HEM+beam & HEM with beam width three. \\
23D feature only & Generic plus role features, without role-submodular
reranking. \\
Structure+greedy & 23D policy with role-submodular completion and greedy
decoding. \\
\method{} & 23D policy, beam width three, policy anchors, and
role-submodular completion. \\
\bottomrule
\end{tabular}
\end{table}

\subsection{Training and implementation}

The frozen training configuration uses five independent training seeds,
60 epochs, Adam learning rates $10^{-4}$ for actor and critic, batch size
eight, maximum gradient norm 2, eight training samples per epoch, and two
parallel sample jobs. The high-level selected-percentage policy uses learning
rate $5\times10^{-4}$. Candidate and selection caps are 128 and 16,
respectively. Checkpoints are selected by validation PDI, not test
performance.

Experiments use SCIP through PySCIPOpt. Main, DOE, SPBS, and sensitivity runs
have a 600-s time limit; OOD runs have a 1200-s limit and a uniform 4096-MB
SCIP memory limit. The final manuscript will report the SCIP, PySCIPOpt,
Python, PyTorch, CUDA, GPU, CPU, operating system, thread count, and the
stage-specific memory limits after all campaigns are completed consistently.
\TBD{Insert the final reproducibility environment table and
repository/archive DOI.}

\subsection{Metrics and statistical analysis}

The primary metric is PDI. Secondary measures are final primal bound, dual
bound, relative gap, primal integral, time to first feasible solution, time
to best incumbent, proof time, root gap, node count, LP iterations, accepted
cuts, callback/inference overhead, timeout rate, memory-failure rate, and
optimal-solution rate.

Comparisons are paired by manufacturing instance and solver seed. Training
seeds characterize learned-model variability but do not convert repeated
solver runs on one manufacturing instance into independent instances. For
each prespecified contrast, we report the paired median difference, a
bootstrap confidence interval over manufacturing instances, the Wilcoxon
signed-rank test where applicable, and Holm-adjusted $p$-values. Timeouts and
failures remain in the dataset and are summarized explicitly; they are never
deleted to improve a method's apparent performance.

\section{Results}
\label{sec:results}

\subsection{Nominal main comparison}

\begin{table}[htbp]
\centering
\caption{Nominal main results. Entries will be computed from 20 instances,
five solver seeds, and, for learned methods, five training seeds. Lower is
better for PDI, gap, and time; higher is better for optimal rate.}
\label{tab:main}
\begin{tabular}{lrrrrr}
\toprule
Method & PDI & Final gap (\%) & First feasible (s) & Nodes & Optimal (\%)\\
\midrule
SCIP & -- & -- & -- & -- & -- \\
ACS & -- & -- & -- & -- & -- \\
HEM & -- & -- & -- & -- & -- \\
HEM+beam & -- & -- & -- & -- & -- \\
23D feature only & -- & -- & -- & -- & -- \\
Structure+greedy & -- & -- & -- & -- & -- \\
\method{} & -- & -- & -- & -- & -- \\
\bottomrule
\end{tabular}
\end{table}

\TBD{Report the complete paired nominal comparison, effect sizes, confidence
intervals, adjusted tests, and failure counts. Do not fill from the currently
partial shards.}

\subsection{Ablation and mechanism analysis}

\begin{table}[htbp]
\centering
\caption{Prespecified mechanism contrasts.}
\label{tab:ablation}
\begin{tabular}{p{0.31\linewidth}p{0.38\linewidth}rr}
\toprule
Contrast & Isolated mechanism & $\Delta$PDI & Holm $p$ \\
\midrule
HEM+beam vs.\ HEM & Search width & -- & -- \\
23D feature only vs.\ HEM+beam & Role features & -- & -- \\
Structure+greedy vs.\ 23D feature only & Submodular completion & -- & -- \\
\method{} vs.\ Structure+greedy & Beam with fixed structural rule & -- & -- \\
\bottomrule
\end{tabular}
\end{table}

\TBD{Add cut-pool size, selected-cut role coverage, redundancy,
callback-overhead, and root-bound analyses to connect solver outcomes to the
proposed mechanism.}

\subsection{Factorial, SPBS, and sensitivity results}

\begin{table}[htbp]
\centering
\caption{Robustness summary over designed manufacturing conditions.}
\label{tab:robustness}
\begin{tabular}{lrrrr}
\toprule
Study & Conditions & Main effect/contrast & $\Delta$PDI & Adjusted $p$\\
\midrule
Factorial DOE & 60 & \TBD{Specify prespecified model term} & -- & -- \\
SPBS & 48 pairs & Automatic vs.\ none & -- & -- \\
Sensitivity & 8 & Worst one-factor degradation & -- & -- \\
\bottomrule
\end{tabular}
\end{table}

\TBD{Report factorial interactions only if prespecified and estimable; report
SPBS separately from the cut-selector comparison because it changes the
manufacturing preprocessing policy.}

\subsection{Out-of-distribution generalization}

\begin{table}[htbp]
\centering
\caption{OOD results by wafer count. All compared methods share the 4096-MB
memory limit and the 1200-s wall-clock limit.}
\label{tab:ood}
\begin{tabular}{lrrrrrr}
\toprule
\multirow{2}{*}{Method} &
\multicolumn{2}{c}{40 wafers} &
\multicolumn{2}{c}{48 wafers} &
\multicolumn{2}{c}{64 wafers} \\
\cmidrule(lr){2-3}\cmidrule(lr){4-5}\cmidrule(lr){6-7}
& PDI & Fail (\%) & PDI & Fail (\%) & PDI & Fail (\%)\\
\midrule
SCIP & -- & -- & -- & -- & -- & -- \\
ACS & -- & -- & -- & -- & -- & -- \\
HEM & -- & -- & -- & -- & -- & -- \\
HEM+beam & -- & -- & -- & -- & -- & -- \\
23D feature only & -- & -- & -- & -- & -- & -- \\
Structure+greedy & -- & -- & -- & -- & -- & -- \\
\method{} & -- & -- & -- & -- & -- & -- \\
\bottomrule
\end{tabular}
\end{table}

At the 27 July 2026 draft audit, 63 of the planned 729 OOD evaluations had
completed under the revised 4096-MB protocol with no recorded run error. This
execution preflight is not used to estimate method performance; the remaining
666 evaluations were running in three memory-budgeted shards. Existing
2048-MB observations are excluded from this campaign. \TBD{After all 729
evaluations finish, report paired outcomes including timeout and
memory-failure incidence.}

\section{Industrial implications}
\label{sec:implications}

The framework separates three responsibilities. Manufacturing engineers
define auditable physical constraints; SCIP retains responsibility for valid
inequalities and proof; the learned component allocates a limited cut budget.
This separation is relevant when an optimizer must be improved without
turning feasibility into a black-box prediction. PDI also reflects the
operational value of obtaining good schedules early, even when a proof of
optimality is not reached within the dispatching horizon.

\TBD{After the complete experiments, translate the observed PDI and
time-to-incumbent improvements into operational implications for the tested
wafer volumes. Avoid extrapolating throughput or cost beyond the modeled
system.}

\section{Limitations and threats to validity}
\label{sec:limitations}

First, the role representation depends on a correct mapping from model
variables to formulation roles. It is portable across instances of the same
model family but has not yet been validated across unrelated MIPs. Second,
the submodular guarantee applies to residual set completion, not solver
runtime. Third, a wall-clock comparison can be hardware- and implementation-
sensitive; hence inference is included in the time limit and the final
environment must be fixed. Fourth, OOD instances create substantial memory
pressure, which must be treated as an outcome rather than censored. Finally,
the study evaluates a simulated deterministic scheduling model; deployment
would additionally require uncertainty, rescheduling, equipment faults, and
fab-level integration.

\section{Conclusions}
\label{sec:conclusion}

This paper presents \method{}, a structure-aware learned cut selector for an
exact dual-source mixed-flow semiconductor scheduling model. The method
combines a 23-dimensional generic-and-role representation, hierarchical
policy anchors, and monotone submodular completion. Its callback operates
only on solver-generated valid cuts and therefore preserves model validity
and exact-solver semantics. A frozen paired protocol has been established to
test nominal performance, mechanisms, designed manufacturing factors, and
distribution shifts. \TBD{Insert the final evidence-based conclusion after
all campaigns and robustness checks are complete.}

\section*{CRediT authorship contribution statement}

\TBD{Insert author-specific CRediT roles.}

\section*{Funding}

\TBD{Insert funding information or state that no specific funding was
received.}

\section*{Declaration of competing interest}

\TBD{Insert the authors' approved declaration.}

\section*{Data and code availability}

The frozen configurations, instance generators, model implementation,
training code, and analysis scripts will be archived at
\TBD{repository URL and immutable DOI}. Raw campaign outputs and a machine-
readable manifest of software/hardware versions will accompany the archive.

\section*{Acknowledgements}

\TBD{Insert acknowledgements or remove this section.}

\bibliographystyle{elsarticle-harv}
\bibliography{rdmct_cie_references}

\end{document}
